\chapter{Αναλυτικά Δεδομένα Λογαριασμών Ρεύματος (\en{Protergia vs NRG})}

Τα δεδομένα που ακολουθούν εξήχθησαν από επίσημους λογαριασμούς ηλεκτρικής ενέργειας (τριφασική παροχή
\texttt{135 kVA}) και χρησιμοποιήθηκαν για την τεχνοοικονομική μελέτη του Κεφαλαίου 6.

\begin{table}[htbp]
\centering
\caption{Περίοδος Βάσης (\en{Before Optimization}) -- Πάροχος \en{Protergia}.}
\label{tab:appendix_energy_baseline}
\begin{tabular}{lccc}
\toprule
\textbf{Μήνας} & \textbf{Ημέρες} & \textbf{Κατανάλωση (\\en{kWh})} & \textbf{Μέση Ημερήσια (\\en{kWh/day})} \\
\midrule
Σεπτέμβριος 2024 & 30 & 4,892 & 163 \\
Οκτώβριος 2024   & 31 & 4,166 & 134 \\
Νοέμβριος 2024   & 30 & 7,552 & 252 \\
Δεκέμβριος 2024  & 31 & 7,869 & 254 \\
Ιανουάριος 2025  & -- & -- & -- \\
Φεβρουάριος 2025 & -- & -- & -- \\
\midrule
\textbf{Σύνολο Περιόδου Βάσης} & \textbf{136} & \textbf{30,309} & \textbf{223} \\
\bottomrule
\end{tabular}
\end{table}

\begin{table}[htbp]
\centering
\caption{Περίοδος Βελτιστοποίησης (\en{After Optimization}) -- Πάροχος \en{NRG}.}
\label{tab:appendix_energy_optimized}
\begin{tabular}{lccc}
\toprule
\textbf{Μήνας} & \textbf{Ημέρες} & \textbf{Κατανάλωση (\\en{kWh})} & \textbf{Μέση Ημερήσια (\\en{kWh/day})} \\
\midrule
Μάρτιος 2025 (Έναρξη)      & 31 & 7,062 & 228 \\
Απρίλιος 2025              & 30 & 5,706 & 190 \\
Μάιος 2025                 & 31 & 3,948 & 127 \\
Ιούνιος 2025               & -- & -- & -- \\
Ιούλιος 2025               & -- & -- & -- \\
Αύγουστος 2025 (Κλειστό)   & 31 & 4,117 & 133 \\
Σεπτέμβριος 2025           & 30 & 5,908 & 197 \\
Οκτώβριος 2025             & 31 & 4,507 & 145 \\
Νοέμβριος 2025             & 28 & 4,749 & 170 \\
\midrule
\textbf{Σύνολο Περιόδου Βελτιστοποίησης} & \textbf{212} & \textbf{37,997} & \textbf{179} \\
\bottomrule
\end{tabular}
\end{table}

\noindent\textit{Σημείωση: Τυχόν κενά ή αποκλίσεις σε επιμέρους μήνες οφείλονται στους ετεροχρονισμένους κύκλους
εκκαθάρισης των παρόχων (\en{Protergia}, \en{NRG}). Τα συνολικά αθροίσματα εξάγονται από τους εκκαθαριστικούς
λογαριασμούς.}

Όπως προκύπτει από την ανάλυση των λογαριασμών, η εφαρμογή του Ψηφιακού Διδύμου επέφερε μείωση της μέσης ημερήσιας
κατανάλωσης από \texttt{223 kWh} σε \texttt{179 kWh}, επιτυγχάνοντας συνολική εξοικονόμηση της τάξης του
\texttt{19.7\%}.
